    \begin{tikzpicture}[>=latex, scale=.8]

        % --- DEFINITION: FLUID ELEMENT MACRO ---
        % This draws a square box with local x, y, z axes.
        % Usage: \fluidelement{x_pos}{y_pos}{rotation_angle}
        \newcommand{\fluidelement}[3]{
            \begin{scope}[shift={(#1,#2)}, rotate=#3]
                % The Box
                \draw[thick, fill=white] (-0.25,-0.25) rectangle (0.25,0.25);
                
                % The Axes (Local coordinates attached to the element)
                % z-axis (up)
                \draw[->] (0,0) -- (0,0.6) node[above, font=\scriptsize] {$z$};
                % y-axis (right)
                \draw[->] (0,0) -- (0.6,0) node[right, font=\scriptsize] {$y$};
                % x-axis (3D projection, bottom-left)
                \draw[->] (0,0) -- (-0.4,-0.4) node[left, font=\scriptsize] {$x$};
            \end{scope}
        }

        % =========================================
        % PART (a): IRROTATIONAL MOTION (Top)
        % =========================================
        
        % 1. The Streamline Path
        \draw[thick, ->] (0, 4) to[out=-10, in=170] (8, 1.5);
        
        % 2. The Fluid Elements (NO ROTATION)
        % Notice the 3rd argument is always "0"
        \fluidelement{1.5}{3.6}{0}
        \fluidelement{3.5}{2.8}{0}
        \fluidelement{5.5}{2.1}{0}
        \fluidelement{7.5}{1.6}{0}
        
        % 3. Label
        \node[anchor=east] at (9, 5) {\textbf{(a) Irrotational frictionless motion}};
        \node[anchor=east, gray] at (9, 4.4) {(Axes maintain orientation)};

        % =========================================
        % PART (b): ROTATIONAL MOTION (Bottom)
        % =========================================
        
        % Shift the whole second diagram down by 5 units
        \begin{scope}[shift={(0,-1)}]
            
            % 1. The Streamline Path
            \draw[thick, ->] (0, 0) to[out=-10, in=170] (8, -2.5);
            
            % 2. The Fluid Elements (WITH ROTATION)
            % The 3rd argument changes (-20, -15, -10, 0) to mimic rotation
            \fluidelement{1.5}{-0.4}{-25}
            \fluidelement{3.5}{-1.2}{-15}
            \fluidelement{5.5}{-1.9}{-5}
            \fluidelement{7.5}{-2.4}{0}
            
            % 3. Label
            \node[anchor=east] at (9, 1) {\textbf{(b) Rotational frictionless motion}};
            \node[anchor=east, gray] at (9, 0.4) {(Axes rotate with the element)};
            
        \end{scope}

    \end{tikzpicture}
